% §11 Exercises
%
% Cross-reference resolution status:
%
% DEFINED here:
%   \label{ex:ch11-1} through \label{ex:ch11-11}
%
% RESOLVED (used here, defined elsewhere):
%   \Cref{sec:miller-rabin-formal}                -- §11.1 (Plan 11 Task 2)
%   \Cref{thm:miller-rabin-formal}                -- §11.1 (Plan 11 Task 2)
%   \Cref{sec:las-vegas-vs-monte-carlo}           -- §11.2 (Plan 11 Task 3)
%   \Cref{def:las-vegas}                          -- §11.2 (Plan 11 Task 3)
%   \Cref{def:monte-carlo}                        -- §11.2 (Plan 11 Task 3)
%   \Cref{sec:rp-bpp}                             -- §11.3 (Plan 11 Task 4)
%   \Cref{def:rp}                                 -- §11.3 (Plan 11 Task 4)
%   \Cref{def:corp}                               -- §11.3 (Plan 11 Task 4)
%   \Cref{def:zpp}                                -- §11.3 (Plan 11 Task 4)
%   \Cref{def:bpp}                                -- §11.3 (Plan 11 Task 4)
%   \Cref{thm:rp-amplification}                   -- §11.3 (Plan 11 Task 4)
%   \Cref{thm:bpp-amplification}                  -- §11.3 (Plan 11 Task 4)
%   \Cref{sec:sample-complexity}                  -- §11.4 (Plan 11 Task 5)
%   \Cref{thm:chernoff}                           -- §11.4 (Plan 11 Task 5)
%   \Cref{sec:miller-rabin}                       -- ch3 §3.4 (Plan 3)
%   \Cref{def:ppv}                                -- ch7 §7.4 (Plan 7)
%   \Cref{def:one-sided}                          -- ch5 §5.4 (Plan 5)
%
% Difficulty-marker macros (not in alex.sty; defined locally if not yet defined):
\providecommand{\exstar}{$\star$}
\providecommand{\exstarstar}{$\star\star$}
\providecommand{\exstarstarstar}{$\star\star\star$}

\section*{Exercises}
\addcontentsline{toc}{section}{Exercises}

\begin{enumerate}

% -----------------------------------------------------------------------
% WARM-UP  (★ × 3): exercises 1, 2, 3
% -----------------------------------------------------------------------

\item \exstar
\label{ex:ch11-1}
\textbf{Miller-Rabin rounds for target FPR.}
\Cref{thm:miller-rabin-formal} bounds the per-round false-positive
rate of Miller-Rabin by $1/4$, so $k$ independent rounds give
$\fprate \le (1/4)^k$. Find the minimum integer $k$ that drives
this bound below $10^{-9}$. Solve $(1/4)^k \le 10^{-9}$ by taking
$\log_{10}$ of both sides: $k \ge 9 / \log_{10} 4 \approx 9 /
0.602 \approx 14.95$, so $k = 15$. Report the value and verify
$(1/4)^{15} \approx 9.3 \times 10^{-10}$.

\item \exstar
\label{ex:ch11-2}
\textbf{Las Vegas vs Monte Carlo identification.}
Classify each of the three algorithms below as Las Vegas (random
runtime, exact output, per \Cref{def:las-vegas}) or Monte Carlo
(fixed runtime, random output, per \Cref{def:monte-carlo}); in
each case, state in one sentence which resource (time or
correctness) is the random quantity.
\begin{enumerate}
  \item[(a)] Randomized quicksort with random pivots.
  \item[(b)] Karger's min-cut algorithm.
  \item[(c)] Miller-Rabin primality test with $k$ fixed rounds.
\end{enumerate}

\item \exstar
\label{ex:ch11-3}
\textbf{RP/BPP/ZPP membership.}
For each language below, identify the smallest complexity class
from $\{\mathrm{P}, \mathrm{ZPP}, \mathrm{RP}, \mathrm{coRP},
\mathrm{BPP}\}$ that contains it, using the definitions of
\Cref{def:rp}, \Cref{def:corp}, \Cref{def:zpp}, \Cref{def:bpp}.
Justify in one sentence each.
\begin{enumerate}
  \item[(a)] \textsc{Compositeness} (input $n$; accept iff $n$ is
    composite), via Miller-Rabin.
  \item[(b)] \textsc{Primes} (input $n$; accept iff $n$ is prime).
  \item[(c)] \textsc{BipartiteMatchingExists} (input bipartite
    graph $G$; accept iff $G$ has a perfect matching).
\end{enumerate}

% -----------------------------------------------------------------------
% MEDIUM  (★★ × 5): exercises 4, 5, 6, 7, 8
% -----------------------------------------------------------------------

\item \exstarstar
\label{ex:ch11-4}
\textbf{Notebook-tied: Miller-Rabin formal vs informal.}
Open the chapter~11 companion notebook
\texttt{python/notebooks/ch11-randomized-algorithms.ipynb}.
\begin{enumerate}
  \item[(a)] Run the formal Miller-Rabin implementation
    (\Cref{sec:miller-rabin-formal}) and the informal walkthrough
    of \Cref{sec:miller-rabin} on $n = 341 = 11 \cdot 31$, a known
    Fermat pseudoprime to base $2$, for $1\,000$ random bases each
    at $k = 1$ round.
  \item[(b)] Report the empirical $\hat{\fprate}$ for each
    implementation and verify both fall below the
    \Cref{thm:miller-rabin-formal} bound of $1/4$.
  \item[(c)] In one sentence, identify the difference between the
    two implementations: the formal version's witness check
    includes the strong-pseudoprime test (composite witnesses
    detected via non-trivial square roots of $1$), while the
    \S3.4 informal version stops at the Fermat test (and so leaks
    on Carmichael numbers in a way the formal version does not).
\end{enumerate}

\item \exstarstar
\label{ex:ch11-5}
\textbf{Chernoff bound application to BPP amplification.}
Apply the \Cref{thm:bpp-amplification} bound to derive the number
of repetitions $k$ needed to reduce a BPP algorithm's error from
$1/3$ to $2^{-10}$.
\begin{enumerate}
  \item[(a)] State the bound: the majority vote of $k$ independent
    runs has error at most $\exp(-k \cdot D(1/2 \| 1/3))$, where
    $D$ is the binary KL divergence (the bound from
    \Cref{thm:bpp-amplification} stated in its sharper form).
  \item[(b)] Compute $D(1/2 \| 1/3) = (1/2) \ln(3/2) +
    (1/2) \ln(3/4) \approx 0.0589$ in natural-log units (nats),
    or use the coarser bound $\exp(-k/18)$ from the standard
    Chernoff form.
  \item[(c)] Solve $\exp(-k/18) \le 2^{-10}$, giving $k \ge 18
    \cdot 10 \cdot \ln 2 \approx 125$. Report $k$ and verify it
    against the bound.
\end{enumerate}

\item \exstarstar
\label{ex:ch11-6}
\textbf{The perfect-filter analog.}
\Cref{def:one-sided} of chapter~5 \S5.4 defines a one-sided
approximate set: $\fnrate = 0$ on the design set $A$, $\fprate
\le \epsilon$ off it. Show that the complexity class $\mathrm{RP}$
of \Cref{def:rp} is the exact analog at the language scale: every
true \textsc{yes}-instance is accepted (witness of class
membership), false-positive accept events occur at most with
probability $1/2$ per run, and amplification by repetition
(\Cref{thm:rp-amplification}) is the operational analog of
combining independent Bloom filters in chapter~6. State in one
sentence what an $\fprate = 0$ analog (perfect filter) would
correspond to: a $\mathrm{coRP}$ language, where one-sided error
on the \textsc{no} side mirrors $\fnrate = 0$ on the design set.

\item \exstarstar
\label{ex:ch11-7}
\textbf{Las Vegas to Monte Carlo conversion via Markov.}
Randomized quicksort is a Las Vegas algorithm with expected
runtime $T_n = O(n \log n)$; the actual mean is $T_n \approx 2 n
\ln n \approx 1.39 n \log_2 n$ comparisons. Apply Markov's
inequality to bound the probability of a long tail.
\begin{enumerate}
  \item[(a)] State Markov: for non-negative $X$ and $a > 0$,
    $\Pr[X \ge a] \le \mathbb{E}[X] / a$.
  \item[(b)] Apply with $X = T_n$ and $a = 10 n \log_2 n$ to
    obtain $\Pr[T_n \ge 10 n \log_2 n] \le 1.39 / 10 = 0.139$.
  \item[(c)] Convert to Monte Carlo by truncation: run quicksort
    for at most $10 n \log_2 n$ comparisons and return whatever
    partial result is in hand. The resulting algorithm has fixed
    runtime $O(n \log n)$ and at most $13.9\%$ chance of incorrect
    output, per \Cref{def:monte-carlo}.
\end{enumerate}

\item \exstarstar
\label{ex:ch11-8}
\textbf{PPV for a BPP decision algorithm.}
A BPP algorithm has both-sided error $\fprate = \fnrate = 1/3$
on every input (the worst case allowed by \Cref{def:bpp}).
The algorithm is run on a workload with prevalence $\pi = 0.1$
(one in ten inputs is a true \textsc{yes}-instance). Compute
$\PPV$ via the formula of \Cref{def:ppv}:
\[
  \PPV = \frac{\pi (1 - \fnrate)}{\pi (1 - \fnrate) + (1 - \pi) \fprate}
       = \frac{0.1 \cdot 2/3}{0.1 \cdot 2/3 + 0.9 \cdot 1/3}
       = \frac{1/15}{1/15 + 3/10}
       \approx 0.182.
\]
Report the value and identify, in one sentence, the lesson: at
low prevalence, even a BPP algorithm meeting its $\fprate$ and
$\fnrate$ bounds delivers low $\PPV$, the chapter~7 \S7.4
prevalence-dependence diagnostic applied at the complexity-class
scale.

% -----------------------------------------------------------------------
% CHALLENGING / EXTENSION  (★★★ × 3): exercises 9, 10, 11
% -----------------------------------------------------------------------

\item \exstarstarstar
\label{ex:ch11-9}
\textbf{$\mathrm{ZPP} = \mathrm{RP} \cap \mathrm{coRP}$.}
Prove that the zero-error probabilistic polynomial-time class of
\Cref{def:zpp} is exactly $\mathrm{RP} \cap \mathrm{coRP}$, using
the Las-Vegas-to-Monte-Carlo conversion lemma of
\Cref{sec:las-vegas-vs-monte-carlo}.
\begin{enumerate}
  \item[(a)] ($\subseteq$) Suppose $L \in \mathrm{ZPP}$, so $L$
    has a Las Vegas decider $M$ with expected polynomial runtime
    and zero error. Truncate $M$ to $2 \cdot \mathbb{E}[T_M(n)]$
    steps and accept by Markov; report \textsc{yes} on
    \textsc{don't-know} to place $L$ in $\mathrm{RP}$, then
    repeat the same truncation reporting \textsc{no} on
    \textsc{don't-know} for $\mathrm{coRP}$.
  \item[(b)] ($\supseteq$) Suppose $L \in \mathrm{RP} \cap
    \mathrm{coRP}$ via deciders $M_+, M_-$. Run $M_+$ and $M_-$
    in alternating rounds; halt when one of them returns its
    certain side (an \textsc{rp} \textsc{yes} or a \textsc{corp}
    \textsc{no}). Show the expected number of rounds is $\le 2$
    by independence, so the combined algorithm is Las Vegas with
    expected polynomial runtime and zero error.
  \item[(c)] Conclude $\mathrm{ZPP} = \mathrm{RP} \cap
    \mathrm{coRP}$, the complexity-class analog of $\fnrate = 0$
    \emph{and} $\fprate = 0$ at the language scale.
\end{enumerate}

\item \exstarstarstar
\label{ex:ch11-10}
\textbf{Karger amplification.}
Karger's contraction algorithm returns the min-cut of an $n$-vertex
graph with probability at least $p_n = 2 / (n (n - 1))$, a Monte
Carlo algorithm in the sense of \Cref{def:monte-carlo}.
\begin{enumerate}
  \item[(a)] Run Karger $r$ independent times and return the
    smallest cut found. The probability that all $r$ runs miss the
    min-cut is at most $(1 - p_n)^r \le \exp(-r p_n) = \exp(-2 r /
    (n (n - 1)))$.
  \item[(b)] Solve $\exp(-2 r / (n (n - 1))) \le \delta$ for $r$,
    giving $r \ge (n (n - 1) / 2) \ln(1 / \delta)$. For $n = 100$
    and $\delta = 0.01$, compute $r \ge 4950 \cdot \ln(100)
    \approx 22\,800$. The cubic dependence on $n$ is the cost of
    amplification when the per-trial success rate is only $1 / n^2$.
  \item[(c)] In one sentence, contrast with Miller-Rabin, where
    per-round $\fprate \le 1/4$ is independent of input size and
    amplification needs only $O(\log(1 / \delta))$ rounds; the
    Karger case is the regime where Monte Carlo amplification is
    expensive because the per-trial signal is weak.
\end{enumerate}

\item \exstarstarstar
\label{ex:ch11-11}
\textbf{Chernoff bound proof.}
Reproduce the moment-generating-function proof of
\Cref{thm:chernoff} (multiplicative Chernoff upper tail). Let
$X_1, \ldots, X_n$ be independent Bernoulli random variables with
$\mathbb{E}[X_i] = p_i$ and $X = \sum_i X_i$, $\mu = \mathbb{E}[X]
= \sum_i p_i$.
\begin{enumerate}
  \item[(a)] For any $t > 0$, apply Markov to $e^{t X}$ to get
    $\Pr[X \ge (1 + \delta) \mu] \le \mathbb{E}[e^{t X}] /
    e^{t (1 + \delta) \mu}$.
  \item[(b)] Use independence and $\mathbb{E}[e^{t X_i}] = 1 +
    p_i (e^t - 1) \le \exp(p_i (e^t - 1))$ (since $1 + x \le
    e^x$) to obtain $\mathbb{E}[e^{t X}] \le \exp(\mu (e^t - 1))$.
  \item[(c)] Optimize $t$ by setting $t = \ln(1 + \delta)$ to
    obtain
    \[
      \Pr[X \ge (1 + \delta) \mu] \le \left(\frac{e^{\delta}}
        {(1 + \delta)^{1 + \delta}}\right)^{\mu},
    \]
    the upper-tail form of \Cref{thm:chernoff}. The lower tail
    follows by the symmetric argument with $-t$ and the
    substitution $\delta \mapsto -\delta$.
\end{enumerate}

\end{enumerate}
